\paragraph{Figure 5.} The proof that $\Mm$ is dense goes as follows.
\begin{enumeratei}
\item[(a)] We first cover $X$ by topologically trivial open sets (here $Y_1,Y_2,Y_3$) on which the degeneracies of $H$ reduce to those of a $2\times2$ system.
\item[(b)] In $Y_1$, the degenerate part of $H$ reduces to that of a $2\times2$ system. Via the procedure of \S3.1, we can produce $H_1$, arbitrarily close to $H$, with no bad points in $Y_1$. By Lemma 3.3, $\BB(H_1)$ is a small perturbation of $\BB(H)\setminus Y_1$.
\item[(c)] We repeat the procedure and produce $H_2$, arbitrarily close to $H_1$, with no bad points in $Y_2$. As bad points are stable, $\BB(H_2)$ is close to $\BB(H_1)$. In particular passing from $H_1$ to $H_2$ does not generate bad points back in $Y_1\setminus(Y_2\cup Y_3)$, and removes bad points in $Y_2$.
\item[(d)] We get new systems $H_1,H_2,H_3$, recursively constructed, arbitrarily close to $H$, with no bad points in $Y_1$, $Y_1\cup Y_2\setminus Y_3$, $Y_1\cup Y_2\cup Y_3$, respectively. Since $Y_1\cup Y_2\cup Y_3$ cover $X$, $H_3$ is in $\Mm$ and is arbitrarily close to $H$.
\end{enumeratei}
